% =====================================================================================================
% appL_lightray.tex  —  附录 L  光线算符与定域算符：矩关系的全阶证明
% writer W09 (spec v2.3.1; LEAD_DECISIONS LA-1…LA-16). Checks: writing/work/W09/v8_appL.py (72/72 PASS),
% v11_appKL_extra.py (b),(f),(g) (22/22 PASS), v4_moments.py; research/c_pdf s4 (8/8 PASS).
% Circularity check of N04 / MD7 (FINAL_OUTLINE_v2 0.12): the proof of prop:appL-renormalized uses only the
% renormalization of LOCAL operators (thm:prelim-composite-renorm); see the paragraph after its proof.
% =====================================================================================================
\chapter{光线算符与定域算符：矩关系的全阶证明}\label{app:appL}

本附录是\cref{thm:pdf-moments-local}的证明：PDF 的 Mellin 矩等于扭度 2 定域算符的矩阵元，对裸算符与对 $\MSbar$ 重整化的量都成立。三步各成一节：顶点层次的均差恒等式（\cref{lem:appL-divided-differences}，纯数学）；裸算符的矩关系（\cref{prop:appL-bare}，$[R,R_{\mathrm{def}}]$）；重整化后的矩关系与与靶无关的卷积核（\cref{prop:appL-renormalized}，$[R,R_{\mathrm{PT}}]$）。证明不预设非定域算符的 BRST 分类（见\cref{sec:appL-renormalized}末）。记号同\cref{ch:pdf}：$w=w^-n_B$，$\partial_{w^-}=n_B\cdot\partial=\partial^+$，$A^+=n_B\cdot A$，$\partial_{w^-}W[w,0]=\ii gA^+(w)W[w,0]$。

% -----------------------------------------------------------------------------------------------------
\section{顶点：均差恒等式}\label{sec:appL-vertices}

设算符顶点不注入动量（前向矩阵元中由整体动量守恒保证）。$m$ 个胶子依次附着在 Wilson 线上，动量 $q_1,\dots,q_m$ 流入（$q_1$ 最靠近 $\psi(0)$），夸克 $k_1$ 由 $\psi(0)$ 流入；记 $\kappa_0=k_1^+$，$\kappa_i=k_1^++q_1^++\dots+q_i^+$，流出夸克的加号动量为 $\kappa_m$；简记 $\kappa_{a:b}\equiv(\kappa_a,\dots,\kappa_b)$。均差递归地定义为 $\phi[\kappa_0]=\phi(\kappa_0)$，$\phi[\kappa_{0:j}]=(\phi[\kappa_{1:j}]-\phi[\kappa_{0:j-1}])/(\kappa_j-\kappa_0)$（重合节点取连续延拓）；互异节点时 $\phi[x_0,\dots,x_k]=\sum_j\phi(x_j)/\prod_{l\ne j}(x_j-x_l)$。$h_r(\kappa_{0:m})=\sum_{0\le i_1\le\dots\le i_r\le m}\kappa_{i_1}\cdots\kappa_{i_r}$ 是完全齐次对称多项式（$h_0=1$，$r<0$ 时 $h_r=0$）。

\begin{lemma}[均差恒等式]\label{lem:appL-divided-differences}
记 $e_w(x)=\ee^{\ii xw}$。(a) $x^p[\kappa_{0:m}]=h_{p-m}(\kappa_{0:m})$（$p\ge0$）。
(b) 对一切实数 $w$，定向的嵌套积分 $E_m(w)\equiv\ii^m\int_0^w\!\dd s_m\int_0^{s_m}\!\dd s_{m-1}\cdots\int_0^{s_2}\!\dd s_1\,\ee^{\ii\kappa_mw-\ii\sum_jq_j^+s_j}$ 等于 $e_w[\kappa_{0:m}]$。
(c) 对固定的附着次序（色因子 $t^{a_m}\cdots t^{a_1}$），$\int\frac{\dd w^-}{4\pi}\ee^{-\ii\xi P^+w^-}O_q(w^-)$ 的 $m$ 胶子顶点为 $g^mn_B^{\mu_m}\cdots n_B^{\mu_1}t^{a_m}\cdots t^{a_1}\gamma^+\cdot\frac12\delta(x-\xi P^+)[\kappa_{0:m}]$（均差作用于 $x$），其 $n$ 阶矩
\begin{equation}\label{eq:appL-vertex}
\int\dd\xi\,\xi^{n-1}(\cdots)=\frac{g^mn_B^{\mu_m}\cdots n_B^{\mu_1}t^{a_m}\cdots t^{a_1}\gamma^+\,h_{n-1-m}(\kappa_{0:m})}{2(P^+)^n}
\end{equation}
恰为 $\bar\psi\gamma^+(\ii D^+)^{n-1}\psi$ 的 $m$ 胶子顶点（同一附着次序）除以 $2(P^+)^n$；$m\ge n$ 时两者都为零。$m=1$ 时 $h_{n-2}(k_1^+,k_2^+)=\sum_{j=0}^{n-2}(k_2^+)^j(k_1^+)^{n-2-j}$，$k_2=k_1+q_1$。
\uses{def:pdf-lightray,prop:appB-rules}
\end{lemma}
% VF-PDF-12 (m = 1). General m: v8 (a) (m <= 4, n <= 9), v11 (b) (induction step, w of both signs, m <= 3).

\begin{proof}
(a) 对 $m$ 归纳（$m=0$ 时两边都是 $\kappa_0^p$）。把 $h_r(\kappa_{0:j})$ 的单项式按是否含 $\kappa_j$ 分类，得 $h_r(\kappa_{0:j})=h_r(\kappa_{0:j-1})+\kappa_jh_{r-1}(\kappa_{0:j})$；按是否含 $\kappa_0$ 分类，得 $h_r(\kappa_{0:j})=h_r(\kappa_{1:j})+\kappa_0h_{r-1}(\kappa_{0:j})$。两式相减得 $h_r(\kappa_{1:j})-h_r(\kappa_{0:j-1})=(\kappa_j-\kappa_0)h_{r-1}(\kappa_{0:j})$，代入均差的递归式即可；重合节点由连续性。
(b) 对 $m$ 归纳（$m=0$ 时两边都是 $\ee^{\ii\kappa_0w}$）。分出最外层积分并用 $\kappa_m-q_m^+=\kappa_{m-1}$，得 $E_m(w)=\ii\int_0^w\ee^{\ii\kappa_m(w-s)}E_{m-1}(s)\,\dd s$，$E_{m-1}$ 以 $\kappa_{0:m-1}$ 为节点。对单个节点 $x$，积分 $\ii\int_0^w\ee^{\ii\kappa_m(w-s)}\ee^{\ii xs}\dd s$ 等于 $(\ee^{\ii xw}-\ee^{\ii\kappa_mw})/(x-\kappa_m)=e_w[x,\kappa_m]$（$w$ 取任意符号）。积分对 $x$ 线性，由归纳假设 $E_m(w)=g_w[\kappa_{0:m-1}]$，$g_w(x)=e_w[x,\kappa_m]$；由显式公式与 $\sum_{j=0}^{k}1/\prod_{l\ne j}(x_j-x_l)=0$（常数的 $k\ge1$ 阶均差）可验证 $(\phi[\,\cdot\,,y])[x_0,\dots,x_{k-1}]=\phi[x_0,\dots,x_{k-1},y]$，故 $E_m(w)=e_w[\kappa_{0:m}]$。
(c) 由\cref{prop:appB-rules}所用的路径排序展开（后面的点在左），$W[w,0]$ 的 $m$ 胶子项是 $(\ii g)^m$ 乘以 $A^+(s_mn_B)\cdots A^+(s_1n_B)$ 在 (b) 的定向嵌套积分区域上的积分。代入 $A^+(sn_B)=\int\frac{\dd^4q}{(2\pi)^4}\ee^{-\ii q^+s}\tilde A^+(q)$ 与 $\bar\psi(w)$ 的因子 $\ee^{\ii\kappa_mw^-}$，$w^-$ 依赖为 $g^mE_m(w^-)=g^me_{w^-}[\kappa_{0:m}]$。均差与 Fourier 变换可交换，而 $\int\frac{\dd w^-}{4\pi}\ee^{-\ii\xi P^+w^-}\ee^{\ii xw^-}=\frac12\delta(x-\xi P^+)$，这给出顶点；又 $\int\dd\xi\,\xi^{n-1}\frac12\delta(x-\xi P^+)=x^{n-1}/(2(P^+)^n)$，由 (a) 得\cref{eq:appL-vertex}。另一方面 $\ii D^+=\ii\partial^++gA^+$；$(\ii D^+)^{n-1}$ 中含按 $A^+(q_m)\cdots A^+(q_1)$ 排列的 $m$ 个 $gA^+$ 的项，由各段 $\ii\partial^+$ 的个数 $p_0,\dots,p_m$ 标记（$\sum_ip_i=n-1-m$，$p_i$ 个位于 $A^+(q_i)$ 与 $A^+(q_{i+1})$ 之间）；这样的 $\ii\partial^+$ 作用于其右侧的 $\psi(k_1),A^+(q_1),\dots,A^+(q_i)$，给出 $\kappa_i$，故该项为 $g^m\prod_i\kappa_i^{p_i}$，求和得 $g^mh_{n-1-m}(\kappa_{0:m})$。
\end{proof}

$m=1$ 即\cref{lem:pdf-wilson-moments}中的单胶子顶点，$m=2$ 给出二阶均差。胶子光线算符的伴随 Wilson 线给出同样的均差，色矩阵换为 $(T^a)_{bc}=-\ii f^{abc}$；场强 $G^{+j}$ 中的胶子场是两边共有的。

% -----------------------------------------------------------------------------------------------------
\section{裸算符的矩}\label{sec:appL-bare}

\begin{proposition}[裸算符恒等式 \tagR]\label{prop:appL-bare}
（$R_{\mathrm{def}}$：由算符定义直接推出。）
设 $|P\rangle$ 是使裸分布 $f_q$、$f_g$ 支撑在 $[-1,1]$ 上的态。则 $M_q(w^-)$、$M_g(w^-)$（\cref{def:pdf-lightray}）是整函数，并且对 $n\ge1$（胶子 $n\ge2$）
\begin{gather}
\int_{-1}^1\dd\xi\,\xi^{n-1}f_q(\xi)=\frac{1}{2P^+}\Big(\frac{-\ii}{P^+}\Big)^{n-1}M_q^{(n-1)}(0)=\frac{\langle P|\bar\psi\gamma^+(\ii D^+)^{n-1}\psi|P\rangle}{2(P^+)^n},\label{eq:appL-bare-q}\\
\int_{-1}^1\dd\xi\,\xi^{n-1}f_g(\xi)=\frac{(-\ii)^{n-2}M_g^{(n-2)}(0)}{(P^+)^n}=\frac{\langle P|O^{(n)}_g|P\rangle}{(P^+)^n},\label{eq:appL-bare-g}
\end{gather}
其中 $O^{(n)}_g=\sum_jG^{+j}_a[(\ii D^+)^{n-2}G^{+j}]_a$，场、耦合与矩阵元都是裸的，伴随表示中 $(D^+X)_a=\partial^+X_a+gf^{abc}A^+_bX_c$。\cref{eq:appL-bare-q}左边等于 $\int_0^1x^{n-1}[q+(-1)^n\bar q]$；$f_g$ 是奇分布，奇数 $n$ 时\cref{eq:appL-bare-g}两边为零，偶数 $n$ 时 $\int_0^1x^{n-1}g=\langle O^{(n)}_g\rangle/(2(P^+)^n)$。以左作用导数写，\cref{eq:appL-bare-q}中间一项等于 $\langle P|\bar\psi(-\ii\overleftarrow{D}{}^+)^{n-1}\gamma^+\psi|P\rangle/(2(P^+)^n)$；两种写法只在前向矩阵元中相等。
\uses{lem:appL-divided-differences,def:pdf-lightray,def:pdf-quark,def:pdf-gluon,def:dis-twist}
\end{proposition}

\begin{proof}
(1) 由\cref{def:pdf-lightray}，$M_q$ 与 $M_g$ 分别是 $2P^+f_q(\xi)$ 与 $(P^+)^2\xi f_g(\xi)$ 的 Fourier 变换（核 $\ee^{\ii\xi P^+w^-}$）。$f_q$、$\xi f_g$ 支撑在 $[-1,1]$ 上，故 $M_q,M_g$ 是整函数\cite[定理~7.3.1]{Hormander2003}，可在配对号下求导：$M_q^{(n-1)}(0)=2P^+(\ii P^+)^{n-1}\int\xi^{n-1}f_q$，$M_g^{(n-2)}(0)=(P^+)^n\ii^{n-2}\int\xi^{n-1}f_g$，即两式的第一个等号。

(2) 由 $\partial_{w^-}W[w,0]=\ii gA^+(w)W[w,0]$，$\partial_{w^-}[\bar\psi(w)W[w,0]]=[\bar\psi\overleftarrow{D}{}^+](w)W[w,0]$，$\bar\psi\overleftarrow{D}_\mu\equiv\partial_\mu\bar\psi+\ii g\bar\psi A_\mu$；$\bar\psi\overleftarrow{D}{}^+$ 与 $\bar\psi$ 同样变换，可以迭代，在 $w=0$ 处 $W=1$，故 $M_q^{(n-1)}(0)=\langle P|\bar\psi(\overleftarrow{D}{}^+)^{n-1}\gamma^+\psi|P\rangle$（左作用形式；顶点层次即\cref{lem:appL-divided-differences}(c)）。
(3) 对按基础表示变换的 $\chi,\chi'$，$\partial^+(\bar\chi\gamma^+\chi')=(\bar\chi\overleftarrow{D}{}^+)\gamma^+\chi'+\bar\chi\gamma^+(D^+\chi')$。记 $V_{k,l}=\bar\psi(-\ii\overleftarrow{D}{}^+)^k\gamma^+(\ii D^+)^l\psi$，则 $V_{k+1,l}=V_{k,l+1}-\ii\partial^+V_{k,l}$。前向矩阵元中 $\langle P|\partial^+V|P\rangle=0$（平移不变性），故 $\langle V_{n-1,0}\rangle=\langle V_{0,n-1}\rangle$，即第二个等号；非前向时 $\langle P'|\partial^+V|P\rangle=\ii(P'-P)^+\langle P'|V|P\rangle\ne0$。

(4) 伴随 Wilson 线满足 $\partial_{w^-}W^{(A)}_{ab}[w,0]=-gf^{acd}A^+_c(w)W^{(A)}_{db}[w,0]$，故 $G^{+j}_a(w)W^{(A)}_{ab}[w,0]$ 对 $w^-$ 的导数是 $[(D^+G^{+j})_a](w)W^{(A)}_{ab}[w,0]$；迭代得 $M_g^{(n-2)}(0)=\sum_j\langle[(D^+)^{n-2}G^{+j}]_aG^{+j}_a\rangle$。由 $\partial^+(X_aY_a)=(D^+X)_aY_a+X_a(D^+Y)_a$，前向矩阵元中每把一个导数移到右边产生一个负号；代入 (1)，
\[
\int\xi^{n-1}f_g=\frac{(-\ii)^{n-2}}{(P^+)^n}\big\langle[(D^+)^{n-2}G]_aG_a\big\rangle=\frac{(-\ii)^{n-2}(-1)^{n-2}}{(P^+)^n}\big\langle G_a[(D^+)^{n-2}G]_a\big\rangle=\frac{\langle O^{(n)}_g\rangle}{(P^+)^n}.
\]由 $f_{\bar q}(\xi)=-f_q(-\xi)$ 与 $f_g$ 的奇性（\cref{def:pdf-quark}、\cref{def:pdf-gluon}）得其余结论。
\end{proof}
% (1),(4): v8 (c),(d) (SU(2), SU(3) adjoint Wilson line), v11 (g); (3): left/right derivatives, c_pdf s4(a).

对红外正规化的部分子态，支撑的假设由\cref{thm:pdf-support-pt}保证；对强子态，它与“光线算符作为缓增分布存在”一起是\cref{asm:intro-composite}(C3) 的内容。左、右作用形式只在前向矩阵元中相等，因此广义部分子分布的矩含有全导数项（\cref{def:excl-gpd}）。

% -----------------------------------------------------------------------------------------------------
\section{重整化后的矩关系}\label{sec:appL-renormalized}

裸恒等式两边要分别重整化：PDF 用卷积核 $Z$（\cref{def:pdf-quark}），定域算符用矩阵 $Z_O(N)$（\cref{thm:prelim-composite-renorm}）。要证明二者相容且 $Z$ 与靶无关。记 C 偶扇区为 $(\{q_k+\bar q_k\},g)$，C 奇扇区为 $\{q_k-\bar q_k\}$；在 $[-1,1]$ 的表述中前者是 $f_{q_k}$ 的奇部分与奇分布 $f_g$，后者是 $f_{q_k}$ 的偶部分，分别只由偶数与奇数阶矩探测。$Z(N)\equiv\int_0^1z^{N-1}Z(z)\,\dd z$。

\begin{proposition}[重整化后的矩关系 \tagR]\label{prop:appL-renormalized}
（$R_{\mathrm{PT}}$：在重整化微扰论中对红外正规化的部分子态严格。）
在维数正规化与 $\MSbar$ 中逐阶地：(a) 存在唯一的味道矩阵值分布 $Z(z)$（支撑在 $[0,1]$ 上，用 $\as(\mu)$ 表示时为 $\delta(1-z)\mathbb 1$ 加 $\varepsilon$ 的纯极点，在每个扇区内作用），使得对一切红外正规化的部分子态 $f^{\mathrm{ren}}=Z\otimes f^{\mathrm{bare}}$ 有限；$Z$ 由一组部分子靶上的极小减除构造，与靶无关。(b) $Z(N)=Z_O(N)$，即 $Z$ 的 Mellin 矩等于扭度 2 定域算符在物理矩阵元中的重整化矩阵（$[O^{(N)}]_\mu=Z_O(N)O^{(N)}$；偶数 $N$ 在 C 偶扇区，奇数 $N$ 在 C 奇扇区）。(c) \cref{eq:appL-bare-q}、\cref{eq:appL-bare-g}对重整化的量成立。（本命题只涉及部分子态。）
\uses{prop:appL-bare,thm:prelim-composite-renorm,def:prelim-msbar,lem:appK-hausdorff,lem:appK-compact,thm:pdf-support-pt}
\end{proposition}

\begin{proof}
在一个扇区内论证，乘积按味道指标理解；“红外正规化”指红外发散由维数正规化以外的方式正规化，且外线满足\cref{thm:prelim-composite-renorm}的条件。
(1) \emph{构造与唯一性。}取一组红外正规化的在壳部分子靶 $j$（扇区的每个分量一个）。矩阵 $F_{ij}=f^{\mathrm{bare}}_{i/j}=\delta_{ij}\delta(1-\xi)+\order{\as}$ 在卷积代数中可逆。递推 $Z^{(k)}=-\mathrm{PP}[F^{(k)}+\sum_{0<l<k}Z^{(l)}\otimes F^{(k-l)}]$（$\mathrm{PP}$ 为 $\varepsilon$ 的极点部分，上标为 $\as$ 的阶）给出纯极点的 $Z$，使 $Z\otimes F$ 逐阶有限。若纯极点的 $Z'$ 也使 $Z'\otimes F$ 有限，则 $Z'\otimes Z^{-1}=(Z'\otimes F)\otimes(Z\otimes F)^{-1}$ 既有限又是 $\mathbb 1$ 加纯极点（这类级数在乘法与求逆下封闭），故为 $\mathbb 1$。由\cref{thm:pdf-support-pt}与卷积保持支撑，$Z$ 支撑在 $[0,1]$ 上（C 偶扇区中 $Z_{gq}$、$Z_{gg}$ 在 $z\to0$ 有 $1/z$ 行为，但该扇区的分布是奇分布，零阶矩为零，卷积有意义）。
(2) \emph{$Z(N)=Z_O(N)$。}对 $[-1,1]$ 上的分布同样有 $\int\xi^{N-1}(Z\otimes f)\,\dd\xi=Z(N)\int\xi^{N-1}f\,\dd\xi$，故 $Z(N)F(N)$ 有限，$F(N)_{ij}=\int\xi^{N-1}f^{\mathrm{bare}}_{i/j}$。由\cref{prop:appL-bare}，$F(N)_{ij}=\langle j|O^{(N)}_i|j\rangle_{\mathrm{bare}}/(2(p^+)^N)$；由\cref{thm:prelim-composite-renorm}，对这些外态扭度 2 算符的矩阵元在 $\{O_q,O_g\}$ 内重整化（与 BRST 恰当算符、运动方程算符与全导数算符的混合在此不出现），即 $Z_O(N)F(N)$ 有限，$Z_O(N)$ 是 $\mathbb 1$ 加纯极点。于是 $Z_O(N)Z(N)^{-1}=[Z_O(N)F(N)][Z(N)F(N)]^{-1}$ 有限且为 $\mathbb 1$ 加纯极点，$Z_O(N)=Z(N)$。
(3) \emph{对一切部分子态有限。}对任一红外正规化的部分子态令 $R=Z\otimes f^{\mathrm{bare}}$。由 (2) 与\cref{prop:appL-bare}，$\int\xi^{N-1}R=Z_O(N)\int\xi^{N-1}f^{\mathrm{bare}}=\langle[O^{(N)}]_\mu\rangle/(2(P^+)^N)$ 有限。逐阶地，$R$ 的每个极点系数是支撑在 $[-1,1]$ 上（\cref{thm:pdf-support-pt}）、有确定奇偶性的分布，其该奇偶的矩全为零，另一奇偶的矩自动为零，由\cref{lem:appK-compact}（测度情形即\cref{lem:appK-hausdorff}）为零。故 $R$ 有限，(a) 成立且 $Z$ 与靶无关。(4) 由 (2)(3)，取矩与重整化可交换，得 (c)。
\end{proof}

单圈时 (b) 可逐图且对 $\varepsilon$ 精确地检验：梯形图、Wilson 线图、自能图分别对应顶角图、算符—胶子图、自能图，$n$ 阶矩之和的极点为 $\frac{\as}{2\pi}\frac{\gamma_{qq}(n)}{\varepsilon}$（\cref{ex:pdf-moments-exact}；我们对 $n\le14$ 逐一核对）。(2) 用到\cref{thm:prelim-composite-renorm}的条件：外线须红外正规化（或动量转移非零）；对在壳无质量胶子的前向矩阵元，若只用维数正规化同时正规化红外发散，与外来（BRST 恰当）算符的混合不必退耦（Hamberg 与 van Neerven 的难题\cite{Hamberg:1991qt}，及\cite{Collins:1994ee}）。对强子态，同样的代数在\cref{asm:intro-wightman}与\cref{asm:intro-composite}(C3) 之下成立，即\cref{thm:pdf-moments-local}的非微扰地位行。
% VF-PDF-13; v8 (e).

\paragraph{非定域算符的分类没有被预设。}
证明只用到\emph{定域}规范不变算符的重整化（\cref{thm:prelim-composite-renorm}）、裸恒等式、支撑与矩的唯一性；与靶无关的核 $Z$ 由部分子靶构造，其重整化性质经矩从定域算符“继承”而来，并未假设光线算符在 BRST 意义下乘法重整化。所以\cref{ch:front}以本命题为输入作矩约化（\cref{prop:front-brst-bilocal}）时不构成循环。光线算符是定域算符的生成函数、重整化的光线算符由重整化的定域算符给出，这一思路在一般的非前向情形见\cite[式~(2.4)、(2.6)、(3.1)、(3.2)]{Braun:2019qtp}；\cite[第~11.4~节]{Collins2011}对 PDF 的重整化断言了同样的三角混合形式。

\cref{thm:pdf-moments-local}的证明至此完成：裸部分是\cref{prop:appL-bare}（支撑由\cref{thm:pdf-support-pt}保证），重整化部分与“与靶无关的核 $Z$ 存在且 $Z(N)=Z_O(N)$”是\cref{prop:appL-renormalized}。
